\documentclass[aspectratio=169,11pt]{beamer}

\usetheme{Madrid}
\usecolortheme{default}
\usefonttheme{professionalfonts}
\setbeamertemplate{navigation symbols}{}
\setbeamertemplate{footline}[frame number]

\usepackage{amsmath, amssymb, booktabs, graphicx}

\newcommand{\E}{\mathbb{E}}
\newcommand{\1}{\mathbf{1}}

\title{Regression Discontinuity, RKD, Bunching, and Local Designs}
\subtitle{Empirical Methods --- Week 5}
\author{Teaching deck draft}
\date{}

\begin{document}

\begin{frame}
  \titlepage
\end{frame}

\begin{frame}{Week 5 question}
\begin{itemize}
    \item What can be learned when rules, thresholds, kinks, dates, or borders create local variation?
    \item What exactly is discontinuous: treatment levels, treatment probabilities, incentive slopes, or densities?
    \item Why should potential outcomes be smooth at the threshold absent treatment?
    \item How local is the estimate, and is that local margin economically meaningful?
\end{itemize}

\medskip
RD-style methods are design arguments before they are software commands.
\end{frame}

\begin{frame}{Local designs as a family}
\small
\begin{tabular}{p{0.22\linewidth} p{0.30\linewidth} p{0.33\linewidth}}
\toprule
Design & Identifying object & Core threat \\
\midrule
Sharp RD & jump in outcome at cutoff & sorting or non-smooth potential outcomes \\
Fuzzy RD & local Wald ratio & weak or opaque first-stage jump \\
RKD & slope discontinuity ratio & noisy derivative estimation, bunching \\
Bunching & excess mass near kink or notch & model dependence, counterfactual density \\
Time RD & jump around a date & seasonality, serial correlation, anticipation \\
Spatial RD & jump across a border & geographic sorting, boundary confounds \\
\bottomrule
\end{tabular}
\end{frame}

\begin{frame}{Sharp vs fuzzy RD}
\[
D_i=\1\{X_i\ge c\}
\]

\[
\tau_{SRD}
=
\lim_{x\downarrow c}\E[Y_i\mid X_i=x]
-
\lim_{x\uparrow c}\E[Y_i\mid X_i=x]
\]

\[
\tau_{FRD}
=
\frac{
\lim_{x\downarrow c}\E[Y_i\mid X_i=x]
-
\lim_{x\uparrow c}\E[Y_i\mid X_i=x]
}{
\lim_{x\downarrow c}\E[D_i\mid X_i=x]
-
\lim_{x\uparrow c}\E[D_i\mid X_i=x]
}
\]

\begin{itemize}
    \item Sharp RD: cutoff fully determines treatment.
    \item Fuzzy RD: cutoff shifts treatment probability; estimate is local to threshold compliers.
\end{itemize}
\end{frame}

\begin{frame}{What RD identifies}
\begin{itemize}
    \item Identifies a local effect at the cutoff under continuity and no precise manipulation.
    \item Does not identify effects far from the cutoff without a portability argument.
    \item Does not identify mechanisms if crossing the threshold changes several channels.
    \item Does not survive sorting, heaping, or manipulation that changes composition at the cutoff.
    \item Strong applied interpretation names the threshold, local population, treatment margin, and unsupported claims.
\end{itemize}
\end{frame}

\begin{frame}{Local polynomial estimation, bandwidths, and weights}
\small
Let $R_i=X_i-c$. A local linear RD solves a weighted least-squares problem:
\[
\min_{\alpha_-,\alpha_+,\beta_-,\beta_+}
\sum_i K\!\left(\frac{R_i}{h}\right)\1\{|R_i|\le h\}
\left[
Y_i-\alpha_- -\beta_-R_i
-\1\{R_i\ge0\}\left(\alpha_+-\alpha_-+(\beta_+-\beta_-)R_i\right)
\right]^2
\]

\[
\widehat{\tau}_{RD}=\widehat{\alpha}_+-\widehat{\alpha}_-
\]

\begin{itemize}
    \item Bandwidth $h$: larger means more precision and more bias risk.
    \item Triangular kernel: $K(u)=(1-|u|)\1\{|u|\le1\}$, emphasizing observations nearest the cutoff.
    \item Report bandwidth, kernel, polynomial order, and effective observations on each side.
\end{itemize}
\end{frame}

\begin{frame}{\texttt{rdrobust} and \texttt{rdhonest}}
\begin{columns}[T,onlytextwidth]
\begin{column}{0.48\linewidth}
\textbf{\texttt{rdrobust} logic}
\begin{itemize}
    \item Estimate local-polynomial RD.
    \item Estimate leading bias with a higher-order fit.
    \item Re-center the estimate and adjust standard errors.
    \item Report conventional and robust bias-corrected inference.
\end{itemize}
\end{column}
\begin{column}{0.48\linewidth}
\textbf{\texttt{rdhonest} logic}
\begin{itemize}
    \item State a smoothness or curvature class.
    \item Bound worst-case bias inside the bandwidth.
    \item Build confidence intervals valid under that class.
    \item Often wider, but assumptions are explicit.
\end{itemize}
\end{column}
\end{columns}

\medskip
They address related but different inference questions.
\end{frame}

\begin{frame}{Covariates and polynomial pitfalls}
\begin{itemize}
    \item Predetermined covariates can improve precision; they do not create RD identification.
    \item Covariates should be measured before treatment and be continuous at the cutoff.
    \item If covariates change the estimate sharply, treat that as a diagnostic warning.
    \item Avoid post-treatment controls, even if they improve fit.
    \item High-order global polynomials are discouraged because far-away data can shape boundary behavior.
    \item Prefer local linear or local quadratic fits with transparent bandwidth sensitivity.
\end{itemize}
\end{frame}

\begin{frame}{Design diagnostics}
\small
\begin{tabular}{p{0.25\linewidth} p{0.58\linewidth}}
\toprule
Diagnostic & What it asks \\
\midrule
Institutional rule & Is the threshold real, binding, and measured correctly? \\
Running-variable density & Is there sorting, heaping, or strategic timing near the cutoff? \\
Covariate continuity & Are predetermined characteristics smooth through the threshold? \\
Bandwidth sensitivity & Does the interpretation depend on one narrow window? \\
Placebo cutoffs & Are similar jumps appearing where no rule changes? \\
Inference level & Are mass points, clusters, time dependence, or spatial dependence handled? \\
\bottomrule
\end{tabular}
\end{frame}

\begin{frame}{Time RD vs event study}
\begin{columns}[T,onlytextwidth]
\begin{column}{0.48\linewidth}
\textbf{Time RD}
\begin{itemize}
    \item Running variable is calendar time.
    \item Cutoff is the treatment date.
    \item Identification is local continuity around the date.
    \item Best with sharp, unanticipated, isolated policy timing.
\end{itemize}
\end{column}
\begin{column}{0.48\linewidth}
\textbf{Event study / DID}
\begin{itemize}
    \item Uses comparison paths over time.
    \item Identification is untreated trend credibility.
    \item Best when untreated units or donor paths are strong.
    \item Handles broader dynamics when comparison paths are valid.
\end{itemize}
\end{column}
\end{columns}

\medskip
Time RD must confront seasonality, serial correlation, anticipation, and concurrent shocks.
\end{frame}

\begin{frame}{Spatial RD}
\[
\tau_{border}
=
\lim_{d\downarrow 0}\E[Y_i\mid \text{signed distance}_i=d]
-
\lim_{d\uparrow 0}\E[Y_i\mid \text{signed distance}_i=d]
\]

\begin{itemize}
    \item One-dimensional version: signed distance to a border.
    \item Multidimensional version: local comparison around a boundary curve or segments.
    \item Credibility requires local geographic continuity and limited sorting.
    \item Implementation often needs boundary-segment fixed effects, local windows, and spatial dependence corrections.
    \item The estimate is local to the border, not automatically a regional or national ATE.
\end{itemize}
\end{frame}

\begin{frame}{RKD and bunching}
\small
\[
\tau_{RKD}
=
\frac{
\lim_{x\downarrow c}\frac{d}{dx}\E[Y_i\mid X_i=x]
-
\lim_{x\uparrow c}\frac{d}{dx}\E[Y_i\mid X_i=x]
}{
\lim_{x\downarrow c}\frac{d}{dx}\E[B_i\mid X_i=x]
-
\lim_{x\uparrow c}\frac{d}{dx}\E[B_i\mid X_i=x]
}
\]

\[
\widehat B=\sum_{z\in \mathcal B}\left[f(z)-\widehat f_0(z)\right],
\qquad
\widehat b=\frac{\widehat B}{\widehat f_0(z^*)}
\]

\begin{itemize}
    \item RKD uses a slope change in incentives or treatment intensity.
    \item Bunching uses excess mass near a kink or notch plus an optimization model.
    \item Kinks change marginal incentives; notches create a level discontinuity and often a dominated region.
\end{itemize}
\end{frame}

\begin{frame}{Choosing among local designs}
\small
\begin{tabular}{p{0.22\linewidth} p{0.31\linewidth} p{0.32\linewidth}}
\toprule
Use when & Preferred design & Key caveat \\
\midrule
Treatment jumps at cutoff & sharp RD & local external validity \\
Eligibility shifts take-up & fuzzy RD & first-stage content \\
Marginal incentive slope changes & RKD & derivative estimation and smoothness \\
Choices pile up at a kink/notch & bunching & behavioral model and counterfactual density \\
Policy begins on exact date & time RD & seasonality, anticipation, serial correlation \\
Institution changes at border & spatial RD & border confounds and sorting \\
\bottomrule
\end{tabular}
\end{frame}

\begin{frame}{Block 1 causal-inference summary}
\small
\begin{tabular}{p{0.23\linewidth} p{0.38\linewidth} p{0.25\linewidth}}
\toprule
Method & Counterfactual logic & Main vulnerability \\
\midrule
Selection on observables & conditionally similar units & unobserved selection, hidden spillovers \\
Experiments & random assignment & noncompliance, attrition, interference \\
DID / event study / SCM & untreated paths or donor pools & trend failures, anticipation, spillovers \\
IV / 2SLS & instrument-induced treatment variation & exclusion, weak first stage, local compliers \\
RD / RKD / bunching & local threshold or density comparison & sorting, locality, implementation sensitivity \\
\bottomrule
\end{tabular}

\medskip
Any credibility ranking is heuristic and context-dependent. Design fit beats method labels.
\end{frame}

\begin{frame}{Research Lab design}
\begin{columns}[T,onlytextwidth]
\begin{column}{0.32\linewidth}
\textbf{Reproduce}
\begin{itemize}
    \item Close-election synthetic RD
    \item Local linear estimate
    \item Triangular weights
    \item RD plot
\end{itemize}
\end{column}
\begin{column}{0.32\linewidth}
\textbf{Diagnose}
\begin{itemize}
    \item Bandwidth sensitivity
    \item Covariate continuity
    \item Density checks
    \item RBC-style and honest intervals
\end{itemize}
\end{column}
\begin{column}{0.32\linewidth}
\textbf{Transfer}
\begin{itemize}
    \item Tax-kink bunching
    \item Counterfactual density
    \item Excess mass
    \item Elasticity memo
\end{itemize}
\end{column}
\end{columns}
\end{frame}

\begin{frame}{Takeaways}
\begin{itemize}
    \item RD-style designs are powerful because the comparison is local and institutional.
    \item Locality must be interpreted, not hidden.
    \item Bandwidths, weights, covariates, polynomial order, and inference define the effective comparison.
    \item Time and spatial RD need extra diagnostics because dates and borders carry their own confounds.
    \item Bunching is a density design plus a behavioral model.
    \item The best threshold papers make clear what the design identifies and what remains outside the threshold.
\end{itemize}
\end{frame}

\end{document}
